% ============================================================
%  CHAPTER 22: CAPACITANCE
%  The Physics Codex: A Complete University Foundation
%  Korede Samuel Omotosho (Falcon)
%  Aurikrex Learning Systems, 2026
%
%  File: chapters/part4/ch22_capacitance.tex
% ============================================================

\chapter{Capacitance}
\label{ch:capacitance}
\minitoc
\newpage

\chapterquote{The art of doing mathematics consists in
finding that special case which contains all the germs
of generality.}{David Hilbert}

\begin{objectivebox}
By the end of this chapter, you should be able to:
\begin{enumerate}[noitemsep]
  \item Define capacitance and the farad
  \item Derive and apply the formula for the
  capacitance of a parallel plate capacitor
  \item Calculate the energy stored in a capacitor
  \item Combine capacitors in series and parallel
  \item Describe the effect of a dielectric on
  capacitance
  \item Describe charging and discharging of a
  capacitor through a resistor
  \item Apply the exponential decay equations for
  capacitor discharge
  \item Describe practical applications of capacitors
\end{enumerate}
\end{objectivebox}

% ============================================================
\section{Intuition: Storing Electric Charge}
% ============================================================

A torch battery stores chemical energy and releases it
as electrical energy. But batteries take time to
recharge and cannot release energy very quickly. A
\textbf{capacitor} stores energy in a completely
different way --- as electric charge on conducting
plates separated by an insulator --- and can release
that energy almost instantaneously.

The camera flash in your phone uses a capacitor.
During the seconds between flashes, the battery slowly
charges the capacitor. When you take a photograph,
the capacitor discharges in a fraction of a millisecond,
releasing a burst of energy far more powerful than the
battery could sustain continuously. The power is
borrowed from the future --- stored over seconds,
released in microseconds.

Capacitors are in every electronic device ever made.
Understanding them is not optional for anyone who wants
to understand electronics.

% ============================================================
\section{Definition of Capacitance}
% ============================================================

\begin{definitionbox}[Capacitance]
The \textbf{capacitance} $C$ of a conductor (or pair
of conductors) is the ratio of the charge stored to
the potential difference applied:
\[
  C = \frac{Q}{V}
  \implies Q = CV
\]
SI unit: farad (F) where $1\ \text{F}
= 1\ \text{C\,V}^{-1}$.

Practical capacitors range from picofarads (pF)
to millifarads (mF):
\[
  1\ \text{pF} = 10^{-12}\ \text{F},\quad
  1\ \text{nF} = 10^{-9}\ \text{F},\quad
  1\ \mu\text{F} = 10^{-6}\ \text{F},\quad
  1\ \text{mF} = 10^{-3}\ \text{F}
\]

Capacitance depends only on the \textit{geometry}
of the conductors and the material between them ---
not on the charge or voltage.
\end{definitionbox}

\begin{historybox}[The Leyden Jar --- First Capacitor]
The first capacitor was invented in 1745 by Ewald
von Kleist and independently by Pieter van Musschenbroek
in Leiden, Netherlands. It consisted of a glass jar
coated inside and out with metal foil. A metal rod
through the stopper connected to the inner foil.
When charged from an electrostatic machine and then
touched, it gave a violent shock. Benjamin Franklin
used Leyden jars in his famous kite experiment.
The Leyden jar is a simple parallel plate capacitor
with the glass acting as the dielectric.
\end{historybox}

% ============================================================
\section{The Parallel Plate Capacitor}
% ============================================================

\begin{processbox}[Capacitance of a Parallel Plate Capacitor]
Two parallel conducting plates, each of area $A$,
separated by distance $d$, with vacuum (or air)
between them.

The surface charge density on each plate:
$\sigma = Q/A$

Electric field between the plates (from Chapter 21):
$E = \sigma/\varepsilon_0 = Q/(A\varepsilon_0)$

Potential difference across the gap:
$V = Ed = Qd/(A\varepsilon_0)$

Capacitance:
\[
  C = \frac{Q}{V} = \frac{Q}{Qd/(A\varepsilon_0)}
\]
\[
  \boxed{C = \frac{\varepsilon_0 A}{d}}
\]
$C$ increases with larger area $A$ (more charge
stored per volt) and decreases with larger separation
$d$ (weaker field per unit distance, less charge
per volt).
\end{processbox}

\begin{diagbox}[Parallel Plate Capacitor: Geometry and Parameters]
\begin{center}
\begin{tikzpicture}[scale=1.0, font=\small]

  % Positive plate (top)
  \fill[codexred!30] (0,3.5) rectangle (7,4.0);
  \draw[very thick] (0,3.5) -- (7,3.5);
  \draw[very thick] (0,4.0) -- (7,4.0);
  \foreach \x in {0.4,1.0,1.6,2.2,2.8,3.4,4.0,4.6,5.2,5.8,6.4} {
    \node[codexred, font=\bfseries] at (\x,3.75) {$+$};
  }
  \node[right, codexred] at (7.1,3.75) {$+Q$};

  % Negative plate (bottom)
  \fill[codexblue!30] (0,0.5) rectangle (7,1.0);
  \draw[very thick] (0,0.5) -- (7,0.5);
  \draw[very thick] (0,1.0) -- (7,1.0);
  \foreach \x in {0.4,1.0,1.6,2.2,2.8,3.4,4.0,4.6,5.2,5.8,6.4} {
    \node[codexblue, font=\bfseries] at (\x,0.75) {$-$};
  }
  \node[right, codexblue] at (7.1,0.75) {$-Q$};

  % Electric field lines (downward, uniform)
  \foreach \x in {0.8,1.8,2.8,3.8,4.8,5.8,6.8} {
    \draw[->, thick, codexgold]
      (\x,3.5) -- (\x,1.0);
  }
  \node[codexgold, font=\small] at (3.5,2.25)
    {$\vect{E} = V/d$};

  % Area label
  \draw[<->, gray, thick] (0,4.3) -- (7,4.3);
  \node[above, gray] at (3.5,4.4) {Area $A$};

  % Separation label
  \draw[<->, gray, thick] (7.8,1.0) -- (7.8,3.5);
  \node[right, gray] at (7.9,2.25) {$d$};

  % Voltage
  \draw[thick, dashed, gray] (-0.5,3.75) -- (0,3.75);
  \draw[thick, dashed, gray] (-0.5,0.75) -- (0,0.75);
  \draw[<->, gray, thick] (-0.5,0.75) -- (-0.5,3.75);
  \node[left, gray] at (-0.6,2.25) {$V$};

  % Formula
  \node[codexblue, font=\large, align=center]
    at (3.5,-0.5)
    {$C = \dfrac{\varepsilon_0 A}{d}$};

  % With dielectric
  \node[font=\small, align=center] at (3.5,-1.3)
    {With dielectric (relative permittivity
    $\varepsilon_r$):
    $C = \dfrac{\varepsilon_r\varepsilon_0 A}{d}$};

\end{tikzpicture}
\end{center}
\end{diagbox}

\begin{examplebox}[22.1]
\stepcounter{workedexample}
A parallel plate capacitor has square plates of side
$15\ \text{cm}$ separated by $2.0\ \text{mm}$ of air.
($\varepsilon_0 = 8.85\times10^{-12}\ \text{F\,m}^{-1}$)\\
(a) Calculate the capacitance.\\
(b) A potential difference of $400\ \text{V}$ is
applied. Find the charge stored.\\
(c) The separation is halved. What is the new
capacitance?\\
(d) What p.d.\ would be needed to store the same
charge as in (b) with the new separation?
\end{examplebox}

\begin{solutionbox}
\textbf{(a)} $A = (0.15)^2 = 0.0225\ \text{m}^2$;
$d = 2.0\times10^{-3}\ \text{m}$:
\[
  C = \frac{\varepsilon_0 A}{d}
  = \frac{8.85\times10^{-12} \times 0.0225}
    {2.0\times10^{-3}}
  = \frac{1.991\times10^{-13}}{2.0\times10^{-3}}
  = 9.95\times10^{-11}\ \text{F}
  \approx 99.5\ \text{pF}
\]

\textbf{(b)} Charge:
\[
  Q = CV = 9.95\times10^{-11} \times 400
  = 3.98\times10^{-8}\ \text{C}
  \approx 40\ \text{nC}
\]

\textbf{(c)} Halving $d$ doubles $C$:
$C' = 2 \times 99.5 = 199\ \text{pF}$

\textbf{(d)} Same charge $Q = 40\ \text{nC}$:
\[
  V' = \frac{Q}{C'} = \frac{3.98\times10^{-8}}
  {1.99\times10^{-10}} = 200\ \text{V}
\]
Halving the separation halves the required voltage ---
the capacitor is more efficient at storing charge.
\end{solutionbox}

% ============================================================
\section{Dielectrics}
% ============================================================

\begin{definitionbox}[Dielectric]
A \textbf{dielectric} is an insulating material placed
between the plates of a capacitor. It increases
capacitance by a factor called the \textbf{relative
permittivity} (dielectric constant) $\varepsilon_r$:
\[
  C = \frac{\varepsilon_r\varepsilon_0 A}{d}
\]
$\varepsilon_r \geq 1$ always.
Air: $\varepsilon_r \approx 1$.
Glass: $\varepsilon_r \approx 5$--$10$.
Ceramic: $\varepsilon_r \approx 10$--$10000$.
Water: $\varepsilon_r \approx 80$.
\end{definitionbox}

\begin{keybox}[Why Dielectrics Increase Capacitance]
When a dielectric is placed in an electric field,
its molecules become \textbf{polarised} --- the
positive and negative charges within each molecule
separate slightly, aligning with the field. This
creates an internal electric field that \textit{opposes}
the applied field. The net field between the plates
is weaker. A weaker field means a lower potential
difference for the same charge. Since $C = Q/V$,
a lower $V$ for the same $Q$ means a \textit{higher}
capacitance.

Dielectrics also allow higher voltages before
breakdown (sparking), enabling more compact capacitors.
\end{keybox}

% ============================================================
\section{Energy Stored in a Capacitor}
% ============================================================

\begin{processbox}[Deriving the Energy Stored in a Capacitor]
When a capacitor is charged from $Q=0$ to $Q=Q_0$,
work must be done against the increasing potential
difference. At intermediate charge $q$, the voltage
is $v = q/C$. The work done adding a small charge
$dq$ is $dW = v\,dq = (q/C)\,dq$.

Total work done (= energy stored):
\[
  W = \int_0^Q \frac{q}{C}\,dq
  = \frac{1}{C}\cdot\frac{Q^2}{2}
  = \frac{Q^2}{2C}
\]
\end{processbox}

\begin{lawbox}[Energy Stored in a Capacitor]
\[
  E = \frac{Q^2}{2C} = \frac{1}{2}QV = \frac{1}{2}CV^2
\]
All three forms are equivalent (use the form that
matches the given information).

The energy is stored in the \textbf{electric field}
between the plates:
\[
  E = \frac{1}{2}\varepsilon_0 E_{field}^2 \times \text{Volume}
  = \frac{1}{2}\varepsilon_0 E_{field}^2 \times Ad
\]
Energy density $= \frac{1}{2}\varepsilon_0 E_{field}^2$
(J\,m$^{-3}$).
\end{lawbox}

\begin{examplebox}[22.2]
\stepcounter{workedexample}
A $100\ \mu\text{F}$ capacitor is charged to
$12\ \text{V}$ by a battery.\\
(a) Find the charge stored.\\
(b) Find the energy stored.\\
(c) The capacitor is disconnected from the battery
and connected to an uncharged $100\ \mu\text{F}$
capacitor. Find the final voltage and total energy.\\
(d) Account for the energy difference between (b)
and (c).
\end{examplebox}

\begin{solutionbox}
\textbf{(a)} $Q = CV = 100\times10^{-6} \times 12
= 1.2\times10^{-3}\ \text{C} = 1.2\ \text{mC}$

\textbf{(b)} Energy stored:
\[
  E = \frac{1}{2}CV^2 = \frac{1}{2}
  \times 100\times10^{-6} \times 144
  = 7.2\times10^{-3}\ \text{J} = 7.2\ \text{mJ}
\]

\textbf{(c)} When connected, charge redistributes.
Total charge is conserved:
$Q_{total} = 1.2\ \text{mC}$.
Total capacitance (parallel): $C_{total} = 200\ \mu\text{F}$.
\[
  V_{final} = \frac{Q_{total}}{C_{total}}
  = \frac{1.2\times10^{-3}}{200\times10^{-6}}
  = 6\ \text{V}
\]
Total energy:
\[
  E_{final} = \frac{1}{2}C_{total}V_{final}^2
  = \frac{1}{2} \times 200\times10^{-6} \times 36
  = 3.6\ \text{mJ}
\]

\textbf{(d)} Energy lost $= 7.2 - 3.6 = 3.6\ \text{mJ}$.
This energy is dissipated as heat in the connecting
wires (even with zero resistance, a spark occurs
during connection which radiates energy). This
result holds regardless of resistance --- exactly
half the initial energy is always lost when two
equal capacitors are connected this way. It is a
consequence of charge conservation combined with
energy minimisation.
\end{solutionbox}

% ============================================================
\section{Capacitors in Series and Parallel}
% ============================================================

\begin{lawbox}[Capacitors in Parallel]
Capacitors in parallel share the same voltage.
The total charge is the sum of charges on each:
\[
  Q_{total} = Q_1 + Q_2 + \ldots
\]
\[
  C_{parallel}V = C_1V + C_2V + \ldots
\]
\[
  \boxed{C_{parallel} = C_1 + C_2 + C_3 + \ldots}
\]
Parallel combination gives a \textbf{larger}
effective capacitance.
\end{lawbox}

\begin{lawbox}[Capacitors in Series]
Capacitors in series carry the same charge.
The total voltage is the sum of voltages across each:
\[
  V_{total} = V_1 + V_2 + \ldots
\]
\[
  \frac{Q}{C_{series}} = \frac{Q}{C_1}
  + \frac{Q}{C_2} + \ldots
\]
\[
  \boxed{\frac{1}{C_{series}}
  = \frac{1}{C_1} + \frac{1}{C_2} + \frac{1}{C_3}
  + \ldots}
\]
Series combination gives a \textbf{smaller}
effective capacitance.
\end{lawbox}

\begin{diagbox}[Capacitors in Series and Parallel]
\begin{center}
\begin{tikzpicture}[scale=0.9, font=\small]

  % === PARALLEL ===
  \begin{scope}[shift={(0,0)}]
    \node[font=\bfseries] at (3.5,4.5)
      {Parallel: $C_p = C_1 + C_2 + C_3$};

    % Bus wires
    \draw[very thick] (0,1.5) -- (0,3.5);
    \draw[very thick] (7,1.5) -- (7,3.5);
    \draw[very thick] (0,3.5) -- (7,3.5);
    \draw[very thick] (0,1.5) -- (7,1.5);

    % Voltage labels
    \node[left] at (-0.2,2.5) {$V$};
    \draw[->, gray] (-0.2,2.5) -- (0,2.5);

    % Three capacitors in parallel
    \foreach \x/\label in {1.2/$C_1$,3.5/$C_2$,5.8/$C_3$} {
      % Wires
      \draw[thick] (\x,3.5) -- (\x,3.0);
      \draw[thick] (\x,2.0) -- (\x,1.5);
      % Plates
      \draw[very thick, codexblue]
        (\x-0.4,3.0) -- (\x+0.4,3.0);
      \draw[very thick, codexred]
        (\x-0.4,2.0) -- (\x+0.4,2.0);
      % Gap
      \node[font=\small] at (\x,2.5) {\label};
    }

    % Same V across each
    \foreach \x in {1.2,3.5,5.8} {
      \node[gray, font=\tiny] at (\x+0.6,2.5) {$V$};
    }
  \end{scope}

  % === SERIES ===
  \begin{scope}[shift={(0,-4)}]
    \node[font=\bfseries] at (5,4.5)
      {Series: $\frac{1}{C_s} = \frac{1}{C_1}
      + \frac{1}{C_2} + \frac{1}{C_3}$};

    % Continuous series path, with one connected wire between each section
    \draw[very thick] (0,2.5) -- (1.75,2.5);
    \draw[thick] (2.25,2.5) -- (4.75,2.5);
    \draw[thick] (5.25,2.5) -- (7.75,2.5);
    \draw[very thick] (8.25,2.5) -- (10.5,2.5);

    % Three capacitors in series; each pair of plates is 0.5 units apart
    \foreach \x/\label in {2/$C_1$,5/$C_2$,8/$C_3$} {
      \draw[very thick, codexblue]
        (\x-0.25,2.1) -- (\x-0.25,2.9);
      \draw[very thick, codexred]
        (\x+0.25,2.1) -- (\x+0.25,2.9);
      \node[above, font=\small] at (\x,2.9) {\label};
    }

    % Voltage labels
    \node[left] at (-0.2,2.5) {$V$};
    \draw[->, gray] (-0.2,2.5) -- (0,2.5);
    \node[below, gray, font=\tiny] at (2,2.1) {$V_1$};
    \node[below, gray, font=\tiny] at (5,2.1) {$V_2$};
    \node[below, gray, font=\tiny] at (8,2.1) {$V_3$};

    % Same Q annotation
    \node[codexgold, font=\small] at (5.25,3.5)
      {Same charge $Q$ on each capacitor};
  \end{scope}

\end{tikzpicture}
\end{center}
\end{diagbox}

\begin{examplebox}[22.3]
\stepcounter{workedexample}
Three capacitors, $C_1 = 4.0\ \mu\text{F}$,
$C_2 = 6.0\ \mu\text{F}$, $C_3 = 12\ \mu\text{F}$,
are connected: $C_1$ in series with the parallel
combination of $C_2$ and $C_3$. A $24\ \text{V}$
battery is connected across the whole arrangement.\\
(a) Find the equivalent capacitance of the network.\\
(b) Find the total charge from the battery.\\
(c) Find the voltage across $C_1$.\\
(d) Find the charge on $C_2$.
\end{examplebox}

\begin{solutionbox}
\textbf{(a)} First: $C_2$ and $C_3$ in parallel:
$C_{23} = 6 + 12 = 18\ \mu\text{F}$

Then $C_1$ in series with $C_{23}$:
\[
  \frac{1}{C_{eq}} = \frac{1}{C_1}
  + \frac{1}{C_{23}} = \frac{1}{4} + \frac{1}{18}
  = \frac{9}{36} + \frac{2}{36} = \frac{11}{36}
\]
\[
  C_{eq} = \frac{36}{11} = 3.27\ \mu\text{F}
\]

\textbf{(b)} Total charge:
\[
  Q_{total} = C_{eq} \times V
  = 3.27\times10^{-6} \times 24
  = 78.5\ \mu\text{C}
\]

\textbf{(c)} In series, $C_1$ carries
$Q = 78.5\ \mu\text{C}$:
\[
  V_1 = \frac{Q}{C_1}
  = \frac{78.5\times10^{-6}}{4\times10^{-6}}
  = 19.6\ \text{V}
\]
Voltage across $C_{23}$: $24 - 19.6 = 4.4\ \text{V}$

\textbf{(d)} Charge on $C_2$ (parallel, so same
voltage $4.4\ \text{V}$ as $C_{23}$):
\[
  Q_2 = C_2 \times V_{23}
  = 6\times10^{-6} \times 4.4
  = 26.4\ \mu\text{C}
\]
\end{solutionbox}

% ============================================================
\section{Charging and Discharging a Capacitor}
% ============================================================

When a capacitor is charged or discharged through a
resistor, the voltage, charge and current do not change
instantly --- they change \textbf{exponentially}.

\begin{processbox}[Deriving the Discharge Equation]
A capacitor of capacitance $C$ is fully charged to
voltage $V_0$ and then discharged through resistance $R$.
At time $t$, let the voltage be $V$.

The discharge current: $I = V/R$ (Ohm's Law)

The rate of charge loss: $I = -dQ/dt = -C\,dV/dt$

(Negative because charge decreases during discharge.)

Setting equal:
\[
  \frac{V}{R} = -C\frac{dV}{dt}
  \implies \frac{dV}{dt} = -\frac{V}{RC}
\]

This is a separable ODE with solution:
\[
  \boxed{V = V_0\,e^{-t/RC}}
\]
Similarly: $Q = Q_0\,e^{-t/RC}$ and
$I = I_0\,e^{-t/RC}$ where $I_0 = V_0/R$.
\end{processbox}

\begin{definitionbox}[Time Constant]
The \textbf{time constant} $\tau = RC$ is the time
for the voltage (or charge) to fall to $1/e
\approx 37\%$ of its initial value.
\[
  \tau = RC \quad [\text{seconds, when }
  R \text{ in ohms and } C \text{ in farads}]
\]
After $\tau$: $V = 0.368\,V_0$ (37\%)\\
After $2\tau$: $V = 0.135\,V_0$ (14\%)\\
After $5\tau$: $V = 0.0067\,V_0$ (0.7\% ---
effectively fully discharged)
\end{definitionbox}

\begin{diagbox}[Charging and Discharging Curves]
\begin{center}
\begin{tikzpicture}
\begin{axis}[
  name=plot1,
  width=11.5cm, height=6cm,
  xlabel={Time $t$ (in units of $\tau = RC$)},
  ylabel={Voltage $V$ (fraction of $V_0$)},
  xmin=0, xmax=5,
  ymin=0, ymax=1.1,
  grid=major, grid style={gray!15},
  xtick={0,1,2,3,4,5},
  xticklabels={$0$,$\tau$,$2\tau$,$3\tau$,$4\tau$,$5\tau$},
  ytick={0,0.368,0.632,1.0},
  yticklabels={0,0.37,0.63,$V_0$},
  legend style={at={(0.5,-0.24)}, anchor=north, font=\scriptsize, /tikz/column 2/.style={column sep=4pt}},
]
  % Discharging: V = V0 * exp(-t/tau)
  \addplot[codexred, very thick, domain=0:5,
    samples=100]
    {exp(-x)};
  \addlegendentry{Discharge: $V = V_0e^{-t/\tau}$};

  % Charging: V = V0 * (1 - exp(-t/tau))
  \addplot[codexblue, very thick, domain=0:5,
    samples=100]
    {1 - exp(-x)};
  \addlegendentry{Charge: $V = V_0(1-e^{-t/\tau})$};

  % 37% marker (discharge)
  \addplot[dashed, codexred!50, domain=0:1]{0.368};
  \draw[dashed, codexred!50]
    (axis cs:1,0) -- (axis cs:1,0.368);
  \fill[codexred] (axis cs:1,0.368) circle (3pt);
  \node[right, codexred, font=\tiny]
    at (axis cs:1.05,0.368)
    {37\% at $t=\tau$};

  % 63% marker (charging)
  \addplot[dashed, codexblue!50, domain=0:1]{0.632};
  \draw[dashed, codexblue!50]
    (axis cs:1,0) -- (axis cs:1,0.632);
  \fill[codexblue] (axis cs:1,0.632) circle (3pt);
  \node[right, codexblue, font=\tiny]
    at (axis cs:1.05,0.632)
    {63\% at $t=\tau$};

  % V0 line
  \addplot[black, dashed, domain=0:5]{1.0};
  \node[right, black, font=\tiny]
    at (axis cs:4.5,1.05) {$V_0$ (battery voltage)};

\end{axis}
\end{tikzpicture}
\end{center}
\small During \textbf{charging}: $V = V_0(1-e^{-t/\tau})$;
$Q = Q_0(1-e^{-t/\tau})$.
During \textbf{discharging}: $V = V_0 e^{-t/\tau}$;
$Q = Q_0 e^{-t/\tau}$.
\end{diagbox}

\begin{examplebox}[22.4]
\stepcounter{workedexample}
A $220\ \mu\text{F}$ capacitor is charged to
$9.0\ \text{V}$ and discharged through a
$47\ \text{k}\Omega$ resistor.\\
(a) Calculate the time constant.\\
(b) Find the voltage after $5.0\ \text{s}$.\\
(c) Find the time at which the voltage has fallen
to $3.0\ \text{V}$.\\
(d) Calculate the initial discharge current.
\end{examplebox}

\begin{solutionbox}
$C = 220\times10^{-6}\ \text{F}$,
$R = 47\times10^3\ \Omega$,
$V_0 = 9.0\ \text{V}$

\textbf{(a)} Time constant:
\[
  \tau = RC = 47\times10^3
  \times 220\times10^{-6}
  = 10.34\ \text{s}
\]

\textbf{(b)} Voltage at $t = 5.0\ \text{s}$:
\[
  V = 9.0\,e^{-5.0/10.34}
  = 9.0\,e^{-0.4836}
  = 9.0 \times 0.617
  = 5.55\ \text{V}
\]

\textbf{(c)} When $V = 3.0\ \text{V}$:
\[
  3.0 = 9.0\,e^{-t/10.34}
  \implies e^{-t/10.34} = \frac{1}{3}
  \implies -\frac{t}{10.34} = \ln\!\left(\frac{1}{3}\right)
  = -1.099
\]
\[
  t = 1.099 \times 10.34 = 11.36\ \text{s}
  \approx 11.4\ \text{s}
\]

\textbf{(d)} Initial current:
\[
  I_0 = \frac{V_0}{R}
  = \frac{9.0}{47\times10^3}
  = 1.91\times10^{-4}\ \text{A}
  = 0.191\ \text{mA}
\]
\end{solutionbox}

% ============================================================
\section{Practical Applications of Capacitors}
% ============================================================

\begin{table}[H]
\centering
\caption{Applications of Capacitors}
\label{tab:capacitor_apps}
\begin{tabular}{lll}
\toprule
\textbf{Application} & \textbf{Capacitor role} &
\textbf{Example}\\
\midrule
Camera flash & Store energy, release rapidly &
  Smartphone flash\\
Radio tuning & Select resonant frequency &
  AM/FM radio receiver\\
Smoothing & Filter ripple from rectified AC &
  Power supply units\\
Timing circuits & Control $RC$ time delay &
  Timer switches, oscillators\\
Memory (DRAM) & Store bits as charge &
  Computer RAM\\
Defibrillator & Deliver precise energy pulse &
  Hospital heart restart\\
Touch screens & Detect finger capacitance &
  Smartphones, tablets\\
Motor starting & Provide phase shift to start &
  Air conditioner, fan\\
Signal coupling & Pass AC, block DC &
  Audio amplifiers\\
\bottomrule
\end{tabular}
\end{table}

\begin{realworldbox}[Capacitors in Nigerian Electronics]
Every inverter used for home backup power during
NEPA outages contains large electrolytic capacitors
(typically $1000$--$10000\ \mu\text{F}$) to smooth
the output of the rectified DC. The UPS (Uninterruptible
Power Supply) in offices across Lagos and Abuja uses
a bank of capacitors to handle the transition between
mains power and battery. The capacitors in Nigerian
power inverters must withstand voltage spikes from
the erratic grid --- capacitor failure is one of the
most common causes of inverter breakdown.
\end{realworldbox}

% ============================================================
\section{Chapter Summary}
% ============================================================

\begin{summarybox}
\textbf{Capacitance:} $C = Q/V$ (farads, F).
Stores charge at a given voltage.

\textbf{Parallel plate:}
$C = \varepsilon_0 A/d$ (vacuum);
$C = \varepsilon_r\varepsilon_0 A/d$ (with dielectric)

$C$ increases with larger area, smaller separation,
higher $\varepsilon_r$.

\textbf{Energy stored:}
$E = \frac{1}{2}CV^2 = \frac{1}{2}QV = \frac{Q^2}{2C}$

\textbf{Parallel:} $C_p = C_1 + C_2 + \ldots$
(larger; same voltage)

\textbf{Series:}
$\frac{1}{C_s} = \frac{1}{C_1} + \frac{1}{C_2}
+ \ldots$ (smaller; same charge)

\textbf{Discharge:}
$V = V_0 e^{-t/RC}$;
$Q = Q_0 e^{-t/RC}$;
$I = I_0 e^{-t/RC}$

\textbf{Charging:}
$V = V_0(1-e^{-t/RC})$

\textbf{Time constant:} $\tau = RC$
(time to reach $37\%$ of initial value during
discharge; $63\%$ of final value during charging)

After $5\tau$: effectively fully charged or
discharged.
\end{summarybox}

% ============================================================
\newpage
\begin{exercisebox}[22]

\subsection*{Section A: Multiple Choice}

\textbf{A1.} A capacitor stores $60\ \mu\text{C}$
at $12\ \text{V}$. Its capacitance is:

\mcoption{A}{$5.0\ \mu\text{F}$}
\mcoption{B}{$6.0\ \mu\text{F}$}
\mcoption{C}{$720\ \mu\text{F}$}
\mcoption{D}{$0.2\ \mu\text{F}$}
\ansref{Ch22-A1}

\vspace{0.4cm}

\textbf{A2.} Three $6\ \mu\text{F}$ capacitors are
connected in series. The effective capacitance is:

\mcoption{A}{$18\ \mu\text{F}$}
\mcoption{B}{$6\ \mu\text{F}$}
\mcoption{C}{$3\ \mu\text{F}$}
\mcoption{D}{$2\ \mu\text{F}$}
\ansref{Ch22-A2}

\vspace{0.4cm}

\textbf{A3.} The energy stored in a $50\ \mu\text{F}$
capacitor charged to $200\ \text{V}$ is:

\mcoption{A}{$0.1\ \text{J}$}
\mcoption{B}{$1.0\ \text{J}$}
\mcoption{C}{$10\ \text{J}$}
\mcoption{D}{$100\ \text{J}$}
\ansref{Ch22-A3}

\vspace{0.4cm}

\textbf{A4.} A capacitor is connected to a battery
and fully charged. The battery is then disconnected.
A dielectric is inserted between the plates. The
voltage across the capacitor:

\mcoption{A}{Increases}
\mcoption{B}{Stays the same}
\mcoption{C}{Decreases}
\mcoption{D}{Drops to zero}
\ansref{Ch22-A4}

\vspace{0.4cm}

\textbf{A5.} A $10\ \mu\text{F}$ capacitor is
charged to $5.0\ \text{V}$ and then connected to
an uncharged $40\ \mu\text{F}$ capacitor. The final
common voltage is:

\mcoption{A}{$1.0\ \text{V}$}
\mcoption{B}{$2.5\ \text{V}$}
\mcoption{C}{$4.0\ \text{V}$}
\mcoption{D}{$5.0\ \text{V}$}
\ansref{Ch22-A5}

\vspace{0.4cm}

\textbf{A6.} During the discharge of a capacitor
through a resistor, the current:

\mcoption{A}{Increases exponentially}
\mcoption{B}{Is constant}
\mcoption{C}{Decreases exponentially}
\mcoption{D}{Decreases linearly}
\ansref{Ch22-A6}

\vspace{0.4cm}

\textbf{A7.} The time constant of a circuit with
$C = 470\ \mu\text{F}$ and $R = 10\ \text{k}\Omega$ is:

\mcoption{A}{$0.047\ \text{s}$}
\mcoption{B}{$4.7\ \text{s}$}
\mcoption{C}{$47\ \text{s}$}
\mcoption{D}{$470\ \text{s}$}
\ansref{Ch22-A7}

\vspace{0.4cm}

\textbf{A8.} For a parallel plate capacitor with
a dielectric, doubling the plate area while halving
the plate separation:

\mcoption{A}{Doubles $C$}
\mcoption{B}{Quadruples $C$}
\mcoption{C}{Halves $C$}
\mcoption{D}{Has no effect on $C$}
\ansref{Ch22-A8}

\vspace{0.4cm}

\textbf{A9.} After three time constants ($3\tau$),
a discharging capacitor retains approximately:

\mcoption{A}{$37\%$ of initial voltage}
\mcoption{B}{$14\%$ of initial voltage}
\mcoption{C}{$5\%$ of initial voltage}
\mcoption{D}{$50\%$ of initial voltage}
\ansref{Ch22-A9}

\vspace{0.4cm}

\textbf{A10.} Two capacitors in parallel both have
voltage $V$ across them. Compared to a single
capacitor of capacitance $C$ at the same voltage,
the total energy stored in two capacitors (each
of capacitance $C$) in parallel is:

\mcoption{A}{The same}
\mcoption{B}{Half}
\mcoption{C}{Double}
\mcoption{D}{Four times}
\ansref{Ch22-A10}

\vspace{0.6cm}

\subsection*{Section B: Theory and Calculations}

\textbf{B1.} A parallel plate capacitor has plates
of area $200\ \text{cm}^2$ separated by an air gap
of $1.5\ \text{mm}$.
($\varepsilon_0 = 8.85\times10^{-12}\ \text{F\,m}^{-1}$)\\[4pt]
(a) Calculate the capacitance.\\
(b) A potential difference of $600\ \text{V}$ is
applied. Find the charge and energy stored.\\
(c) A dielectric of relative permittivity $4.5$
is inserted (battery still connected). Find the
new capacitance, charge and energy.\\
(d) Explain why the energy stored increases when
a dielectric is inserted with the battery connected,
even though the voltage stays the same. Where does
the extra energy come from?
\hfill\ansref{Ch22-B1}

\vspace{0.4cm}

\textbf{B2.} Four capacitors are arranged as shown:
$C_1 = 3\ \mu\text{F}$ and $C_2 = 6\ \mu\text{F}$
in parallel, then this combination in series with
$C_3 = 4\ \mu\text{F}$, then this combination in
parallel with $C_4 = 8\ \mu\text{F}$.
A $20\ \text{V}$ supply is connected.\\[4pt]
(a) Find the capacitance of $C_1$ and $C_2$ in
parallel.\\
(b) Find the capacitance of this with $C_3$ in series.\\
(c) Find the total effective capacitance.\\
(d) Find the total charge from the supply.\\
(e) Find the voltage across $C_3$ and hence the
charge on $C_3$.
\hfill\ansref{Ch22-B2}

\vspace{0.4cm}

\textbf{B3.} A $1000\ \mu\text{F}$ capacitor is
charged to $15\ \text{V}$ by a battery and then
discharged through a $5.0\ \text{k}\Omega$ resistor.\\[4pt]
(a) Calculate the initial charge, initial current
and initial energy.\\
(b) Calculate the time constant.\\
(c) How long does it take for the voltage to fall
to $5.0\ \text{V}$?\\
(d) What is the current at $t = 10\ \text{s}$?\\
(e) Sketch, with numerical labels on the axes,
the discharge curve for voltage vs time from
$t = 0$ to $t = 4\tau$.\\
(f) The resistor is replaced with one of
$10\ \text{k}\Omega$. How does this affect:\\
(i) The initial current; (ii) The time constant;
(iii) The total charge delivered to the resistor.
\hfill\ansref{Ch22-B3}

\vspace{0.4cm}

\textbf{B4.} A defibrillator capacitor bank stores
energy to deliver a shock to a patient's heart.
The capacitors are charged to $5000\ \text{V}$.
The required energy is $360\ \text{J}$.\\[4pt]
(a) Calculate the required capacitance.\\
(b) Calculate the charge stored.\\
(c) If the defibrillator discharges through the
patient's chest resistance of $50\ \Omega$ with a
time constant of $5.0\ \text{ms}$, find the
capacitance (using $\tau = RC$) and check against
your answer to (a).\\
(d) The current through the patient peaks at the
moment of discharge. Calculate this peak current.\\
(e) Calculate the power delivered at the instant
of discharge.\\
(f) After $5\tau$, what fraction of the energy
has been delivered? Comment on the design
implication.
\hfill\ansref{Ch22-B4}

\vspace{0.4cm}

\textbf{B5.} A student investigates the discharge
of a capacitor by measuring the voltage at equal
time intervals:

\begin{center}
\begin{tabular}{cc}
\toprule
$t$ (s) & $V$ (V)\\
\midrule
0 & 8.00\\
5 & 5.90\\
10 & 4.35\\
15 & 3.20\\
20 & 2.36\\
25 & 1.74\\
\bottomrule
\end{tabular}
\end{center}

(a) Taking natural logs: show that
$\ln V = \ln V_0 - t/\tau$.
Calculate $\ln V$ for each row.\\
(b) Plot $\ln V$ against $t$ and draw the best
fit straight line.\\
(c) From the gradient of the line, determine
the time constant $\tau$.\\
(d) From the $y$-intercept, determine $V_0$.\\
(e) The resistor used was $R = 33\ \text{k}\Omega$.
Calculate the capacitance. Compare with the labelled
value of $C = 470\ \mu\text{F}$.
\hfill\ansref{Ch22-B5}

\end{exercisebox}